% META: {"id": "1SPE-TRIGO-EX-040", "chapitre": "1SPE-TRIGONOMETRIE", "type_objet": "exercice", "capacites": ["1SPE-TRIGONOMETRIE-C4"], "capacites_codes": ["C4"], "methodes": ["M4"], "parcours": 3, "competences": ["raisonner", "communiquer"], "duree_min": 18, "mode_creation": "ex_nihilo", "sources_inspiration": [], "parametres_sympy": {}, "fichier_tex": "chapitres/1SPE-TRIGONOMETRIE/exercices/1SPE-TRIGO-EX-040.tex", "corrige_tex": "chapitres/1SPE-TRIGONOMETRIE/corriges/1SPE-TRIGO-CO-040.tex", "status": "generated"}
% BEGIN-VERIFY
% from sympy import *
% x = symbols('x', real=True)
% # 2*cos^2(x) - 3*cos(x) + 1 = 0
% u = symbols('u')
% sols_u = solve(2*u**2 - 3*u + 1, u)
% assert sorted(sols_u) == [Rational(1,2), 1]
% # cos(x) = 1 => x = 2k*pi
% # cos(x) = 1/2 => x = pi/3 + 2k*pi or x = -pi/3 + 2k*pi
% assert cos(pi/3) == Rational(1,2)
% END-VERIFY
\begin{exercice}{1SPE-TRIGO-EX-040}{3}{18}
Resoudre dans $\mathbb{R}$ l'equation $2\cos^2 x - 3\cos x + 1 = 0$.

\begin{enumerate}
  \item Poser $u = \cos x$ et resoudre l'equation du second degre $2u^2 - 3u + 1 = 0$.
  \item Pour chaque valeur de $u$, resoudre l'equation $\cos x = u$.
  \item Donner l'ensemble complet des solutions.
\end{enumerate}
\end{exercice}
